%% =====================================================================
%%  B6-T-13-02 -- what B6 contributes to "Weakness as a Mask"
%%  -------------------------------------------------------------------
%%  LAYER A: APPEND-ONLY. Written once at the entry's write, then left
%%  alone. If the source has more to say later, add a bullet -- do not
%%  rewrite. Material found ONLY here goes in the `unique` slot with
%%  @unique set to yes, which makes it undeletable (rule R10).
%% =====================================================================
%% @kind: source
%% @id: B6-T-13-02
%% @book: B6
%% @topic: T-13-02
%% @locator: ch. 9 (playing the victim, vilifying the victim, servant role), pp. 118--121
%% @status: distilled
%% @unique: no
%% @integrated: yes
%% @updated: 2026-09-19

\dossier{B6}{T-13-02}{\bookshort{B6}}{ch. 9 (playing the victim, vilifying the victim, servant role), pp. 118--121}

\begin{distilled}
  \item Playing the victim: portraying oneself as a victim of circumstance or of someone else's behaviour to gain sympathy, evoke compassion, and thereby get something --- covert-aggressors count on targets never suspecting a persona of weakness could be an instrument.
  \item Vilifying the victim runs with it: the aggressor appears to be merely defending himself against the victim's aggression, which puts the target on the defensive twice.
  \item The servant role is the institutional version: portrayed self-sacrifice as a claim on obedience --- the cornerstone, B6 notes, of corrupt ministerial empires.
  \item B6's deepening of the mask: the displayed weakness is aggression in transport --- the question is never whether the persona is meek but what the meekness is collecting.
\end{distilled}
